\section{Tujuan}

\begin{enumerate}
    \item Memahami fenomena fisis pembentukan inti majemuk $^{13}\text{C}^*$ dan formulasi matematis model resonansi hamburan nuklir satu tingkat Breit-Wigner (\textit{single-level Breit-Wigner formula}).
    \item Menurunkan formulasi analitik eksak vektor gradien $\mathbf{g}(\mathbf{p})$ dan matriks kelengkungan Hessian Gauss-Newton $\mathbf{H}(\mathbf{p})$ untuk minimisasi fungsi objektif kuadrat terkecil terbobot ($\chi^2$).
    \item Mengembangkan dan menerapkan algoritma Newton-Raphson multivariat untuk menyelesaikan sistem linier $\mathbf{H}\Delta\mathbf{p} = -\mathbf{g}$ menggunakan pemrograman ilmiah Python dan NumPy.
    \item Mengekstrak nilai estimasi parameter resonansi optimal ($E_r, \Gamma, \sigma_0, \sigma_{bg}$) dari data eksperimen penampang lintang $^{12}\text{C}(n,\text{tot})$ repositori IAEA-NDS EXFOR.
    \item Mengevaluasi tingkat kelayakan statistik model melalui \textit{reduced chi-square} ($\chi^2_{red}$) serta menghitung ketidakpastian parameter dan derajat korelasi silang dari matriks kovariansi asimptotik $\mathbf{C} = 2\mathbf{H}^{-1}$.
    \item Menguji batas kestabilan dan titik kegagalan (\textit{failure modes}) algoritma komputasi terhadap degradasi data melalui simulasi Monte Carlo uji ketahanan derau (\textit{noise stress-test}).
    \item Memvalidasi hasil komputasi energi resonansi $E_r$ terhadap data standar literatur nuklir internasional IAEA.
\end{enumerate}